% generated by R/04_tables.R from table2_repeated_genes.md -- do not edit
% needs \usepackage{booktabs,adjustbox} in the preamble
\begin{table}[htbp]
\centering
\caption{Genes hit by more than one independent mutation}%
\label{tab:parallel}
\small
\begin{adjustbox}{max width=\linewidth}
\begin{tabular}{llrrrll}
\toprule
\textbf{Gene} & \textbf{Locus tag} & \textbf{Length (bp)} & \textbf{Independent events} & \textbf{Expected} & \textbf{P (Poisson)} & \textbf{P (family-wise)} \\
\midrule
\textit{sinR} & BSU\_24610 & 336 & 5 & 0.0024 & $6.23\times10^{-16}$ & $<5\times10^{-5}$ \\
\textit{ywcC} & BSU\_38220 & 672 & 2 & 0.0047 & $1.12\times10^{-5}$ & 0.00855 \\
\bottomrule
\end{tabular}
\end{adjustbox}
\par\smallskip
{\footnotesize\raggedright Endpoint total fraction. Variants in the same gene carried by the same clone and no other are counted once, since they cannot be shown to be independent -- this is what removes \textit{recX}, whose two variants are both in clone 4 alone. Expected is the number of events a gene of that length would collect if all events were scattered across the coding genome at random. The family-wise column is a min-P permutation over all 4,171 CDS, which is the right correction when the genes were nominated by the data. Gene names and lengths are from the RefSeq annotation of NC\_000964.3.\par}
\end{table}
